\section{Konklusjon}
\label{sec:konklusjon}
Problemstillingen var om bilen kan kjøre minst \SI{5}{m} i snitt over ti kjøringer, med under \SI{15}{\percent} spredning og et sideavvik under \SI{10}{\percent} av lengden. Etter den andre kjøretesten, med ti kjøringer, er målene for distanse og pålitelighet nådd, men ikke målene for spredning og sideavvik. Resultatmålene vurderes slik (tabell~\ref{tab:kravstatus}):
\begin{enumerate}
    \item \textbf{Funksjon og pålitelighet.} Bilen kjørte av seg selv og stoppet uten å treffe noe i alle ti forsøkene i den andre kjøretesten. Målet er nådd.
    \item \textbf{Kjøredistanse.} Bilen kjørte \SI{8,8}{m} i snitt over ti kjøringer, og alle var over \SI{5}{m}. Målet er nådd.
    \item \textbf{Repeterbarhet.} Spredningen var \SI{21}{\percent}, over grensen på \SI{15}{\percent}. Målet er ikke nådd.
    \item \textbf{Retningsstabilitet.} Bilen dro til venstre i alle ti forsøkene, og sideavviket var i snitt \SI{24}{\percent} av lengden. Målet er ikke nådd.
    \item \textbf{Rammer og produksjon.} Bilen bruker det utdelte svinghjulet og akselen, trekkes opp med 20 cm snor og holder seg innenfor $140 \times 85 \times 70$ mm i CAD. Alle delene er laget med 3D-printer og utstyret på verkstedet.
\end{enumerate}

Sideavviket skyldes trolig at den printede rammen er litt vridd, slik at høyre bakhjul og venstre forhjul bærer det meste av vekten. Med en stiv drivaksel tar da høyre drivhjul mest av drivkraften. Da lagersetet for høyre drivhjul ble filt for å jevne ut lasten, ble buene omtrent dobbelt så slake og distansen \SI{61}{\percent} lengre, men bilen dro fortsatt til venstre. For å nå R5 ved denne distansen må svingradiusen være over \SI{44}{m}. Det krever trolig en foraksel som kan vippe, eller tre hjul, slik at hjullasten ikke avhenger av hvor rett rammen er. Spredningen kan forklares av at turtallet etter opptrekk for hånd varierer, og kan bare reduseres med et jevnere opptrekk.

For distansen viser simuleringen at energien fordeler seg på flere tapsledd. Omtrent en firedel går tapt til sluring ved nedsetting, og med godt grep bestemmes den andelen av utvekslingen og massen fremfor friksjonen. En like stor del går i lagrene på svinghjulsakselen. Åpne SKF-lagre ga en åttendedel av tapet med 626 2RS med fett og en firedel av tapet med tørre 626 2RS, og med SKF på svinghjulsakselen øker den simulerte distansen fra 9,4 til \SI{12,4}{m}. Opptrekk med snor ga dessuten bare 1700 til \SI{2850}{o/min}, ikke de \SI{4000}{o/min} som ble antatt i konseptfasen, og en mindre trommel ga ikke høyere turtall. I den andre kjøretesten var tre forsøk lengre enn simuleringen gir ved det høyeste turtallet, så den foreløpige rullemotstanden i simuleringen er trolig satt for høyt.

Iterasjonene med prøvebiter fungerte godt for delene som må passe sammen, som lagersetet, firkantnavet og trommelen. Den første kjøretesten avdekket likevel at den printede rammen var vridd, noe som ga skjev lastfordeling på hjulene og ikke kunne oppdages i delprøvene. Svinghjulet ble først mottatt 2.~november 2026. Flere av delene, blant annet trommelen, måtte tilpasses sveisen og akselen slik de faktisk ble, og bilen kunne ikke kjøres før svinghjulet var på plass. Det begrenset tiden til å designe, iterere og teste, og den første kjøretesten kom derfor bare ni dager før innlevering. Erfaringen er at bilen også må prøves som helhet så tidlig som mulig, og ikke bare del for del.
